\subsection{Negative-weight rigidity from qualitative finiteness}
\label{subsec:selmer-negative-weight-rigidity}

The global rigidity needed at the distinguished representation requires
only qualitative finiteness. No estimate for the size of an algebraic
\(K\)-group enters this step.

\begin{proposition}\label{prop:selmer-negative-weight-rigidity}
Let \(k\) be a number field, let \(\ell\) be any prime, and let
\(\epsilon:G_k\to\overline{\mathbf Q}_\ell^\times\) be a continuous
finite-order character. For every integer \(m\ge1\),
\[
 H_f^1\bigl(k,\overline{\mathbf Q}_\ell(-m)\otimes\epsilon\bigr)=0.
\]
Here the local condition is unramifiedness away from \(\ell\) and the
Bloch--Kato finite condition at \(\ell\), with algebraic coefficient
extensions. There is no restriction on
\(\operatorname{ord}(\epsilon)\) relative to \(\ell\).
\end{proposition}

\begin{proof}
First put a continuous cocycle \(z\) over one finite coefficient field.
Its character has values in a finite extension \(E_0/\mathbf Q_\ell\).
There are finitely many extensions of each degree up to isomorphism,
and finitely many embeddings of each into the algebraic closure.
Thus \(\overline{\mathbf Q}_\ell\) is covered by countably many
finite extensions \(E_i\).\cite[Proposition 7.64 and Remark 7.65]{milne-ant}
The preimages
\[
                 H_i=\{g:z(g)\in E_0E_i\}
\]
are closed subgroups, by the cocycle equation. They cover the compact
Galois group. Baire's theorem gives one with nonempty interior, hence
finite index. Adjoining the cocycle values on its finitely many coset
representatives puts all values in one finite extension \(E\).
This uses the cocycle equation, not a claim about arbitrary compact
subsets of the algebraic closure. Enlarge \(E\) further to contain the
finitely many coefficients of the local crystalline-period witnesses.
The defining coboundary equations then hold over
\(E\otimes_{\mathbf Q_\ell}B_{\mathrm{cris}}\), since this tensor
product injects into the one over the algebraic closure.

Choose a finite Galois extension \(L/k\) killing \(\epsilon\) and
containing \(i\), so \(L\) is totally imaginary. Restriction preserves
both local finite conditions. The restricted character is
\(\theta=\chi_\ell^{-m}\). At a place above \(\ell\), the
Bloch--Kato formula gives
\[
 \dim_{\mathbf Q_\ell} H_f^1(L_v,V)
 =\dim_{\mathbf Q_\ell}\bigl(D_{\mathrm{dR}}(V)/
                                      \operatorname{Fil}^0\bigr)
       +\dim_{\mathbf Q_\ell}H^0(L_v,V).
\]
Both terms vanish for \(V=E(-m)\), so the class is a full-local
coboundary.\cite[Corollary 3.8.4 and Example 3.9]{bloch-kato-tamagawa}
Away from \(\ell\), the coefficient action on inertia is trivial;
an inertia coboundary is therefore zero. Thus the cocycle factors
through \(G_{L,S}\), where \(S\) consists of the \(\ell\)-adic and
archimedean places.

Let \(O\) be the ring of integers of \(E\). Proposition
\ref{shared-selmer-finite}, applied to a \(\mathbf Z_\ell\)-basis
of \(O\), makes
\[
 \mathcal S_E=\ker\left(
 H^1(G_{L,S},(E/O)(-m))
 \longrightarrow\prod_{v\mid\ell}H^1(L_v,(E/O)(-m))\right)
\]
finite. Choose \(a\ge0\) with \(\ell^a\mathcal S_E=0\).
No bound for \(a\) is needed. The cyclotomic character has infinite
image over \(L\), so choose \(g_0\) with
\(d=\theta(g_0)-1\ne0\). Subtract the coboundary with primitive
\(z(g_0)/d\), making \(z(g_0)=0\).
For every \(n\ge a\), the cocycle \(\ell^{-n}z\bmod O\) gives
a class in \(\mathcal S_E\). Hence some \(t_n\in E/O\) satisfies
\[
 \ell^{a-n}z(g)=(\theta(g)-1)t_n\quad\text{in }E/O.
\]
At \(g_0\) this says \(dt_n=0\), so \(t_n\) has a representative
in \(d^{-1}O\). Multiplying by \(d\) gives
\[
                      dz(g)\in\ell^{n-a}O
                              \qquad(n\ge a).
\]
The intersection is zero, so the normalized cocycle vanishes.
No compatibility among the primitives \(t_n\) is required.

The original class is therefore a coboundary on \(G_L\).
Subtract its primitive globally, so \(z|_{G_L}=0\). Normality gives,
for \(h\in G_L\) and \(g\in G_k\),
\[
  \theta(h)z(g)=z(hg)=z\bigl(g(g^{-1}hg)\bigr)=z(g).
\]
One \(h\) has \(\theta(h)\ne1\), so \(z=0\) everywhere. This last
step explains why the order of the finite character and the extension
degree may both be divisible by \(\ell\).
\end{proof}

For the distinguished three-dimensional representation, write
\(\Psi=\chi_\ell^e\psi\), with \(e\ge1\), \(\psi\) finite order,
and \(\eta\) the imaginary quadratic character. With weights
\((1,0,-1)\), put \(R(X)_{ab}=X_{4-b,4-a}\). On weight \(j\),
the adjoint action is \(\Psi(g)^jX\) if \(\eta(g)=1\), and
\(-\Psi(g)^jR(X)\) otherwise. The projections
\((1\pm R)/2\) commute with the action and preserve the finite
local condition. The negative part is therefore exactly
\[
\begin{array}{c|c|c}
 \text{weight}&\text{generator}&\text{character}\\ \hline
 -1&E_{21}-E_{32}&\chi_\ell^{-e}\psi^{-1}\\
 -1&E_{21}+E_{32}&\chi_\ell^{-e}\psi^{-1}\eta\\
 -2&E_{31}&\chi_\ell^{-2e}\psi^{-2}\eta .
\end{array}
\]
Proposition \ref{prop:selmer-negative-weight-rigidity} treats all
three summands. In weight zero, restriction to the imaginary quadratic
field makes the action trivial. The same local calculation identifies
the finite classes with unramified classes at \(\ell\), so Proposition
\ref{prop:selmer-zero-weight-rigidity} supplies that case.

\paragraph{Passing from equations to representation tangents.}
The following comparison supplies point lifting and the lifting and
determination of tangents once the prescribed weighted presentation is available.
All coefficient extensions in its proof are finite and may depend on
the point and tangent vector.

\begin{proposition}\label{prop:selmer-tangent-comparison}
Let \(\mathcal O\) be the ring of integers of a finite extension of
\(\mathbb Q_\ell\), and put \(P=\mathcal O[[z_1,\ldots,z_d]]\),
with integer weights on the variables. Let \(J\) be an ideal stable
under weight components, \(R=P/J\), and \(A\subset R\) the
algebra of finite sums of homogeneous series. Let \(Q\subset A\)
be a finitely generated homogeneous ideal and \(S=R/QR\). Suppose
\((A/Q)_0\) is finite over \(\mathcal O\).

Every \(\overline{\mathbb Q}_\ell\)-point of \(A/Q\) becomes,
after one grading rescaling, the restriction of a point
\(f:S\to\overline{\mathbb Q}_\ell\) satisfying
\(|f(\mathfrak m_S)|<1\). Every derivation on \(A\) at \(f\)
annihilating \(Q\) extends uniquely to a continuous derivation on
\(S\). If \(S\) represents the prescribed framed deformation
problem, its tangent cocycle determines this derivation.
\end{proposition}

\begin{proof}
First let \(B\) be any complete Noetherian local \(\mathcal O\)-algebra
with the same residue field, and \(h:B\to\overline{\mathbb Q}_\ell\)
such a local point. Put \(\mathfrak p=\ker h\). We have
\(|h(B)|\le1\), so \(B/\mathfrak p\) is integral over
\(\mathcal O\). Every prime containing \((\pi)+\mathfrak p^2\)
is therefore maximal, giving
\(\mathfrak m_B^N\subseteq(\pi)+\mathfrak p^2\).
Consequently \(B/\mathfrak p^2\) is \(\pi\)-adically complete
and finite modulo \(\pi\). Successively lifting generators modulo
\(\pi\) makes it finite over \(\mathcal O\). A derivation
\(\delta\) at \(h\) kills \(\mathfrak p^2\), hence has bounded
image in one finite coefficient extension; moreover
\[
 |\delta(\mathfrak m_B^{Nn})|
       \le |\pi|^n\sup_{b\in B}|\delta(b)|\longrightarrow0.
\]
Thus every algebraic derivation at this point is continuous.

Homogeneous series satisfy
\(A=A_0[z_1,\ldots,z_d]\): the exponent vectors of a fixed
weight have finitely many coordinatewise minimal elements, and every
other exponent differs from one of them by a nonnegative vector of
weight zero. Also \(A_0\) is local: a weight-zero series with unit
constant term has a weight-zero inverse. Thus \(B=(A/Q)_0\) is
finite local. Since \(\mathfrak m_B^n\subseteq\pi B\) for some
\(n\), every point \(\xi\) has \(v(\xi(\mathfrak m_B))>0\).
Put \(a_i=\xi(z_i)\) and \(w_i=\operatorname{wt}(z_i)\), omitting zero values. Mixed invariant
monomials give
\[
 \max_{w_i>0}\frac{-v(a_i)}{w_i}
 <r<\min_{w_j<0}\frac{-v(a_j)}{w_j}.
\]
Absent constraints leave an unbounded endpoint; weight-zero coordinates
already have positive valuation. Choose rational \(r\) in this
interval and \(t\) with \(v(t)=r\). Then all \(t^{w_i}a_i\)
have positive valuation. The homogeneous equations are preserved.
Their degree-zero image and finitely many coordinates lie in one finite
extension. The substitution below with no nilpotent part therefore
extends precisely this rescaled point to a local point of \(S\).
With no nonzero coordinate constraints take \(t=1\).

For a tangent at such a point, set
\(\tau(a)=f(a)+\epsilon D(a)\), where \(\epsilon^2=0\).
The finite module \(\tau(A_0)\), together with the
finitely many \(\tau(z_i)\), therefore lies in
\[
 C_b=\mathcal O_E+\epsilon\pi_E^{-b}\mathcal O_E
\]
for some finite extension \(E\) and integer \(b\ge0\).
Its scalar coordinates are integral. The \(\tau(z_i)\) generate
a topologically nilpotent ideal: their scalar coordinates have
absolute value less than one and a product contains at most one
\(\epsilon\)-term. Substitution consequently defines \(P\to C_b\).

This substitution agrees with \(\tau\) on a homogeneous series:
after truncation at total degree \(n\), its tail is a finite sum
of degree-\(n\) monomials times homogeneous series whose images
stay in \(C_b\). The tail tends to zero. Weightwise stability
of \(J\), followed by the same truncation, shows that substitution
kills \(J\). It also kills \(Q\), so descends to \(S\).
Its scalar component equals \(f\) by boundedness and polynomial
density. Its second component is the desired extension. The preceding
continuity argument and density prove uniqueness.

Finally, universal matrix entries topologically generate a framed
deformation ring. Otherwise a nonzero residual tangent functional
annihilating them would give two different maps to the residual dual
numbers with the same representation. Choose finitely many entries
spanning that cotangent space; Nakayama and completeness give the
topological generation. The cocycle is
\(c_\delta(g)=\delta(\rho(g))f(\rho(g))^{-1}\).
If it vanishes, \(\delta\) kills every entry and hence their dense
polynomial algebra. Continuity gives \(\delta=0\).
\end{proof}

\begin{lemma}\label{lem:selmer-normalized-coboundary}
At every prescribed marked point, the primitives of normalized
coboundaries lie in the line spanned by
\(T=\operatorname{diag}(1,0,-1)\).
\end{lemma}

\begin{proof}
Write \(R(X)=J_3X^tJ_3\), with \(J_3\) antidiagonal, and
\(M=\rho(g_1)\). The marking has \(\rho(c)=I\) in the
negative component, \(g_1\) in the positive component, and
\(g_1^*=cg_1^{-1}c\); thus \(\rho(g_1^*)=R(M)\).
Its two marked entries give
\[
 M-R(M)=\zeta T,\qquad \zeta=M_{11}-M_{33}\ne0.
\]
For a coboundary with primitive \(X\), the differentiated marking is
\(X+R(X)=0\) and \([X,\zeta T]_{12}=[X,\zeta T]_{21}=0\).
The first identity gives
\(X=\left(\begin{smallmatrix}a&b&0\\d&0&-b\\0&-d&-a\end{smallmatrix}\right)\);
the last two give \(-\zeta b=\zeta d=0\). Hence \(X=aT\).
\end{proof}

\begin{corollary}\label{lem:selmer-normalized-base-coboundary}
At the base representation, every normalized coboundary is zero.
Under the stated local finite condition and entry-weight compatibility, the nonpositive equation
classes span the corresponding rational cotangent spaces. At any
graded point whose local rescaling has only coboundary cocycles, the tangent
space of the marked equation algebra has dimension at most one.
\end{corollary}

\begin{proof}
The base adjoint action fixes \(T\). Apply the preceding lemma,
Proposition~\ref{prop:selmer-tangent-comparison}, and the nonpositive
vanishing already proved. Entry-weight compatibility preserves the
weight of a base tangent cocycle; cotangent duality gives spanning. At other
rigid points the faithful cocycle map has image in the line of
coboundaries with primitive in \(\overline{\mathbb Q}_\ell T\).
Grading rescaling transports a tangent by \(D_t(a)=t^jD(a)\)
on weight \(j\), preserving its dimension, the equations, and
all degree-zero trace values.
\end{proof}

The weighted presentation, finite degree-zero quotient,
local finite condition, and non-diagonal cohomological rigidity
remain arithmetic inputs. These comparisons give no uniform bound on
integral torsion. The characteristic-zero argument includes
\(\ell=2\); the separate odd-prime restriction of the arithmetic
model is retained.

